\begin{abstract}
\noindent Diese Projektdokumentation untersucht die Anwendbarkeit von tabellarischem Q-Learning im Blackjack-Szenario. Das Hauptaugenmerk liegt auf dem Vergleich zweier unterschiedlicher Architekturen: Einem \textit{Baseline-Agenten}, welcher ausschließlich auf der Standard-Zustandsrepräsentation operiert, und einem erweiterten \textit{Counting-Agenten}, welcher Kartenzählinformationen (\textit{Running Count}, \textit{True Count}) nativ in seinen Zustandsraum integriert. Die Agenten wurden über $100.000.000$ Episoden hinweg unter Verwendung von drei unterschiedlichen Seeds trainiert und anschließend evaluiert. Die Ergebnisse zeigen deutliche Unterschiede in der erlernten Aktionsverteilung sowie der resultierenden Win Rate.
\end{abstract}